\section{Complexity}
\label{sec:complexity}

The decision procedures search for holes up to the bound of
\cref{thm:small-hole}. This section estimates the bound, in the size of the
reader and in the size of the sentence, and compares it with the known lower
bounds. Recall that $\expo_0(x)=x$ and $\expo_{h+1}(x)=2^{\expo_h(x)}$.

\subsection{An Elementary Upper Bound}\label{sec:complexity-upper}

\begin{theorem}[Elementary bounds]\label{thm:elementary}
For all $k\ge0$ and $N\ge1$,
\begin{align*}
  \FreshB(k,N)&\le\expo_2\bigl((k+N+16)^4\bigr), &
  \TailB(k,N)&\le\expo_2\bigl((k+N+20)^{2k+5}\bigr), \\
  \Hole(k,N)&\le\expo_4\bigl((k+N+20)^{2k+6}\bigr).
\end{align*}
The bounds do not depend on the size of the alphabet.
\end{theorem}

\begin{proof}
The inequalities are proved in \cref{app:complexity-closed}, by estimating
the recurrences of \cref{thm:fresh-bound,thm:tails-bound} and the
composition \eqref{eq:hole-bound}.
\end{proof}

In outline: for $\FreshB$, the recurrence $\nu_{d+1}\le 4E F^2\nu_d^2$ of
\cref{thm:fresh-bound} squares its value $kN$ times, so $\log_2\nu_{kN}$ is
at most $2^{kN}$ times a polynomial in $k$ and $N$, and $\FreshB=kX^{kN+1}$
is doubly exponential. For $\TailB$, one shows by induction on the scale
that the logarithms of all scale functions of scale $l$ are bounded by a
common value $H_l$ with $H_{l+1}\le2^{\ell}H_l$, where
$\ell=\ell(k,N)\le(k+N+10)(\Qc+N+3)^k+5$: by \cref{lem:closed-free} the search
length $\Lb_k$ is bounded by a polynomial in the window of degree
$2^{(k+N+10)(\Qc+N+3)^k}$, because the chain budget is the constant $\Qc+N$, and the
other scale functions are polynomials of bounded degree in $\Lb_k$, $N!$ and
the functions of the previous two scales. After the $3(2k+1)+1$ scales, the
logarithms of the period bound $D$ and of $E_0$ are at most
$2^{O(k\ell)}$. Each compression pass replaces $E$ by a value that is linear
in $E$, with a coefficient of the form $2^{k(N-1)}\cdot\mathrm{poly}(D,N!)$
(\cref{thm:tails-bound}), so the $k(\Qc+3)$ passes multiply the logarithm by a
factor polynomial in $k$ and $N$. Finally,
$\Hole(k,N)=\FreshB(k,k(N+\TailB(k,N)+1)+1)$ applies the bound on fresh
tuples to a table whose size is about $k\TailB(k,N)$.

The innermost argument $(k+N+20)^{2k+5}$ is a polynomial in $N$ for each
fixed $k$, but not a polynomial in $k$ and $N$ jointly. So, for each fixed
number $k$ of starts, the bounds on $\FreshB$ and $\TailB$ are doubly
exponential in $N$ and the bound on $\Hole$ is fourfold exponential in $N$.
Uniformly in $k$ and $N$, since $(k+N+20)^{2k+6}=2^{(2k+6)\log_2(k+N+20)}$,
the estimates give $\TailB(k,N)\le\expo_3(O(k\log(k+N+20)))$ and
$\Hole(k,N)\le\expo_5(O(k\log(k+N+20)))$, while $\FreshB(k,N)\le
\expo_2((k+N+16)^4)$ is doubly exponential in a polynomial in $k$ and $N$.
These are upper bounds; we do not know how far they are from the truth.

Three counts in the proof of \cref{thm:tails} keep the bound elementary;
each replaces a cruder count that would give a tower of exponentials whose
height grows with $k$ or $N$.
\begin{itemize}
\item \emph{Suffix nodes instead of windows.} A closed factor is found by a
  pigeonhole on suffix nodes (\cref{def:node}), of which there are
  polynomially many in the window $B$. A pigeonhole on the last $B$ letters
  before a cut would cost $\len\Alp^{B}$ at each of the nested uses in
  \cref{lem:closed-free}, and would make the bound depend on the alphabet.
\item \emph{A constant chain budget.} The suffix chains of a fixed word have
  at most $\Qc+N$ steps from the initial window of the scale on
  (\cref{lem:chains}). The obvious bound on the number of steps, by the
  length of the non-periodic part of the tail, grows from scale to scale,
  and the number of iterations in \cref{lem:closed-free} with it.
\item \emph{A linear cap.} In a compression pass, the simple-path bound
  $2^h$ is applied only to the pieces whose number of pairwise comparable
  members is bounded by states; the pieces in the rejected tails are counted
  linearly (\cref{lem:cap}). Applying the bound $2^h$ to all pieces would
  add one exponential in each of the $k(\Qc+3)$ passes.
\end{itemize}

\begin{corollary}[Sentences]\label{cor:complexity}
Let $\varphi$ be a Boolean combination of $\Sigma_1$ and $\Pi_1$ sentences
of size $m$. Every clause of its normal forms that is realizable over
$\Trees$ is realized by a solution whose witness paths have total length at
most $\expo_5(p(m))$ for a fixed polynomial $p$, and the truth of $\varphi$
over each of the three classes is decidable in space $\expo_5(p(m))$.
\end{corollary}

\begin{proof}
A clause has at most $m$ subterms and at most $m$ negated inequalities, so
its flat constraint has $v\le m$ variables, $v\le m+1$ with the guard of
\cref{sec:decision-bf}, and $k\le m$ queries. The readers have at most
$2^{4v}+1$ states (\cref{thm:reduction}), and
$\expo_4((m+2^{4m+4}+21)^{2m+6})\le\expo_5((4m+6)(2m+6))$, because
$m+2^{4m+4}+21\le2^{4m+6}$. The search of
\cref{thm:decide-arbitrary} enumerates the clauses, the mode vectors and the
tuples one after the other and stores one tuple at a time; the tests of
satisfiability, of finite satisfiability and of the hole conditions need
space polynomial in the size of the reader and the tuple.
\end{proof}

The bound is far from practical. It is a bound on a brute-force search, obtained from
shortening arguments that were not designed for efficiency.

\subsection{Lower Bounds}\label{sec:complexity-lower}

\begin{proposition}[Hardness]\label{prop:hardness}
For a covariant constructor of every fixed arity $n\ge2$, the $\Sigma_1$
theory over $\BFT$ is PSPACE-hard over each of the three classes of trees,
already for plain entailment.
\end{proposition}

\begin{proof}
Plain entailment is a $\Pi_1$ sentence (\cref{sec:prelim-constraints}), and
it is PSPACE-hard: for $n=2$ over infinite trees by Henglein and Rehof
\cite{HengleinRehof98}, and for every $n\ge2$ and each of the three classes
by \cite{Zhao26}, by a reduction from the universality of nondeterministic
finite automata.
\end{proof}

For the arities $0$ and $1$, to which \cref{prop:hardness} does not apply,
Boolean combinations of sentences are still NP-hard, since propositional
satisfiability reduces to them as soon as two distinct trees, such as $\Bot$
and $\Top$, are definable; for $n\ge2$ this is subsumed by the proposition.

\subparagraph*{Long witnesses.}
Shortest holes can be long. For one query and a covariant constructor of
arity $n\ge2$, there are constraints with $O(\beta)$ variables that have
finite solutions and whose shortest witness paths have length exponential in
$\beta$, over each of the three classes \cite{Zhao26}: the constraint lets
an automaton accept every word except the extensions of the record of a
binary counter. With two queries there is a second, different mechanism,
which uses the channels between the two words.

\begin{example}[Witnesses forced by the least common multiple]
\label{ex:lcm}
For $r\ge1$ moduli $p_1,\dots,p_r\ge1$ with least common multiple $L$,
there is a
structural reader with two starts and $O(p_1+\dots+p_r)$ states, over the
alphabet $\{1,2\}$, whose holes are exactly the tuples
\[
  W_t=\bigl((1^{tL}2)^r,\;1^{tL}2\bigr),\qquad t\ge1,
\]
so its shortest hole has total length $(r+1)(L+1)$. The first word passes
$r$ stages; stage $j$ counts a block of letters $1$ modulo $p_j$ and lets
the word proceed only if the block length is divisible by $p_j$. The
admitted pairs are the pairs of a state at the beginning of a block and a
state inside a block, in either order and in both words, which forces all
blocks to have the same length. For a binary covariant constructor, there is
a constraint with $O(p_1+\dots+p_r)$ variables and literals and two queries
that has finite solutions and whose countermodels, over each of the three
classes, have exactly these tuples as their tuples of witness paths; so the
same lower bound holds for the witnesses of countermodels. With the first
$r$ primes as moduli, $L=2^{\Theta(r\log r)}$ while the constraint has
$O(r^2\log r)$ variables. \Cref{app:complexity-long} gives the construction
and the proofs.
\end{example}

\subsection{The Gap}\label{sec:complexity-gap}

Between PSPACE-hardness and the upper bound of \cref{cor:complexity} there
are five exponentials. Already for plain entailment the exact complexity is
open, between PSPACE and 2-EXPSPACE \cite{Zhao26}. For disjunctive
entailment we know no witnesses longer than exponential and no hardness
result beyond PSPACE. The upper bound is unlikely to be tight: the two
shortening arguments are composed through the cover of \cref{sec:overview},
which, for a fixed number of queries, applies a doubly exponential bound to
a table of doubly exponential size, and each of the two bounds is itself obtained by pigeonhole arguments
on finite summaries rather than by a search over summaries. For plain and
for existential entailment, sharper analyses that do not enumerate words give
procedures in 2-EXPSPACE \cite{Zhao26,Zhao26E}; whether the shortening of
tuples admits a similar analysis is open (\cref{sec:future}).
